% =============================================================
%  Module Description — Theory of Computation (CS401)
%  DRAFT for a planned module (Semester 4 of the Foundation programme).
%  The code CS401 is a proposal (semester digit + position in the
%  semester); add the term suffix (e.g. -S27) once the term is known and
%  make it equal the course's `code` in src/content/courses/.
%  Compile with: pdflatex cs401-module-description.tex
%  Write "TODO(owner): ..." where information is still missing;
%  the website highlights it.
% =============================================================
\documentclass[11pt,a4paper]{article}

\usepackage[utf8]{inputenc}
\usepackage[T1]{fontenc}
\usepackage{lmodern}
\usepackage[margin=2.2cm]{geometry}
\usepackage{booktabs}
\usepackage{tabularx}
\usepackage{array}
\usepackage{enumitem}
\usepackage{xcolor}
\usepackage{titlesec}
\usepackage{fancyhdr}
\usepackage[hidelinks]{hyperref}

% ---------- Module details (edit these) ----------------------
% \ModuleCode must equal the course's `code` in src/content/courses/.
\newcommand{\ModuleTitle}{Theory of Computation}
\newcommand{\ModuleCode}{CS401}
\newcommand{\Programme}{Foundation programme}
\newcommand{\Level}{Semester 4}
\newcommand{\Semester}{Spring}
\newcommand{\Credits}{TODO(owner)}
\newcommand{\ModuleLeader}{Dr. Ali Khudiyev}
\newcommand{\Email}{ali.khudiyev@wosm.academy}
\newcommand{\OfficeHours}{Monday--Friday, 10:00--17:00 (Baku time, UTC+4)}
\newcommand{\Prerequisites}{M101-F26 Discrete Structures; CS302 Data Structures and Algorithms II}

% ---------- Styling ------------------------------------------
\definecolor{accent}{RGB}{0,70,130}
\titleformat{\section}{\large\bfseries\color{accent}}{\thesection.}{0.5em}{}
\titlespacing*{\section}{0pt}{1.2em}{0.5em}
\setlist{itemsep=2pt, topsep=4pt}
\setlength{\parindent}{0pt}
\setlength{\parskip}{4pt}

\pagestyle{fancy}
\fancyhf{}
\lhead{\small \ModuleCode\ -- \ModuleTitle}
\rhead{\small Module Description}
\cfoot{\small \thepage}
\renewcommand{\headrulewidth}{0.4pt}

% =============================================================
\begin{document}

\begin{center}
  {\LARGE\bfseries\color{accent} \ModuleTitle}\\[4pt]
  {\large Module Description}
\end{center}
\vspace{6pt}

% ---------- Summary table (not shown on the website: the course page
% shows these details from the macros above) --------------------------
\begin{tabularx}{\textwidth}{>{\bfseries}l X >{\bfseries}l X}
\toprule
Module code   & \ModuleCode    & Credits      & \Credits \\
Programme     & \Programme     & Level        & \Level \\
Semester      & \Semester      & Prerequisites & \Prerequisites \\
Module leader & \ModuleLeader  & Email        & \href{mailto:\Email}{\Email} \\
Office hours  & \multicolumn{3}{l}{\OfficeHours} \\
\bottomrule
\end{tabularx}

% ---------- Sections (shown on the website) ------------------
\section{Module Overview}
Theory of Computation asks what computers can do in principle and what they can do
efficiently. It studies abstract models of computation, from finite automata to
Turing machines, and the classes of languages they recognise. Finite automata, met
briefly in Digital Design and Computer Architecture I, are treated here in full.

The module has three parts. \emph{Automata and formal languages} covers regular and
context-free languages, which underlie pattern matching, lexing and parsing.
\emph{Computability} shows that some problems, such as the halting problem, cannot
be solved by any program. \emph{Complexity} makes precise the distinction between
tractable and intractable problems introduced in Data Structures and Algorithms II.

\section{Aims}
\begin{itemize}
  \item Introduce the standard mathematical models of computation and the relationships between them.
  \item Show the limits of computation through rigorous proofs of undecidability.
  \item Give a precise foundation for complexity classes and NP-completeness.
  \item Develop your ability to write careful proofs about formal objects.
\end{itemize}

\section{Learning Outcomes}
After completing this course, you will be able to:
\begin{enumerate}[label=LO\arabic*.]
  \item Construct deterministic and nondeterministic finite automata and regular expressions, convert between them, and minimise automata.
  \item Prove that languages are not regular using the pumping lemma and the Myhill--Nerode theorem.
  \item Construct context-free grammars and pushdown automata, transform grammars into Chomsky normal form, parse with the CYK algorithm, and prove that languages are not context-free.
  \item Describe Turing machines and equivalent models, explain the Church--Turing thesis, and prove problems undecidable by diagonalisation and reduction.
  \item Place problems in complexity classes such as P, NP and PSPACE, and prove that problems are NP-complete by reduction.
\end{enumerate}

\section{Syllabus}
\begin{tabularx}{\textwidth}{c X}
\toprule
\textbf{Unit} & \textbf{Topic} \\
\midrule
1  & Alphabets, words and languages; proof techniques revisited \\
2  & Deterministic finite automata \\
3  & Nondeterminism and the subset construction \\
4  & Regular expressions and Kleene's theorem \\
5  & Non-regular languages: pumping lemma, Myhill--Nerode, minimisation \\
6  & Context-free grammars, derivations and ambiguity \\
7  & Chomsky normal form and the CYK algorithm \\
8  & Pushdown automata and their equivalence with context-free grammars \\
9  & Non-context-free languages; closure properties \\
10 & Turing machines and their variants \\
11 & The Church--Turing thesis; universal Turing machines \\
12 & Decidability and the halting problem \\
13 & Reductions; Rice's theorem; the Post correspondence problem \\
14 & Time complexity and the class P \\
15 & The class NP and the Cook--Levin theorem \\
16 & NP-complete problems and reductions \\
17 & Space complexity: PSPACE and Savitch's theorem \\
18 & Course Review \& Exam \\
\bottomrule
\end{tabularx}

\section{Teaching and Learning Methods}
The module is self-paced. All materials are provided on Moodle.
Each unit has a lecture and exercise sheets with constructions and proofs.

\begin{tabularx}{0.6\textwidth}{X r}
\toprule
\textbf{Activity} & \textbf{Hours} \\
\midrule
Lectures            & 36 \\
Labs / Tutorials    & 24 \\
Independent study   & 140 \\
\midrule
\textbf{Total}      & \textbf{200} \\
\bottomrule
\end{tabularx}

\section{Assessment}
Your knowledge will be assessed throughout the course with quizzes, exercise sheets
and the final exam. The exam checks whether you can build automata and grammars for
given languages, and whether you can write correct proofs of non-regularity,
undecidability and NP-completeness.

\begin{tabularx}{\textwidth}{X c c c}
\toprule
\textbf{Component} & \textbf{Weight} & \textbf{Due} & \textbf{LOs} \\
\midrule
Quizzes & 20\% & Throughout & LO1--LO5 \\
Exercise sheets & 30\% & Throughout & LO1--LO5 \\
Final exam & 50\% & End of module & LO1--LO5 \\
\midrule
\textbf{Total} & \textbf{100\%} & & \\
\bottomrule
\end{tabularx}

\section{Reading List}
\textbf{Essential}
\begin{itemize}
  \item M.\ Sipser: \textit{Introduction to the Theory of Computation}, 3rd ed., Cengage Learning, 2012.
\end{itemize}
\textbf{Recommended}
\begin{itemize}
  \item J.\ E.\ Hopcroft, R.\ Motwani, J.\ D.\ Ullman: \textit{Introduction to Automata Theory, Languages, and Computation}, 3rd ed., Pearson, 2006.
  \item D.\ C.\ Kozen: \textit{Automata and Computability}, Springer, 1997.
  \item S.\ Arora, B.\ Barak: \textit{Computational Complexity: A Modern Approach}, Cambridge University Press, 2009.
\end{itemize}

\section{Materials and Licence}
Course materials are published under the
\href{https://creativecommons.org/licenses/by-sa/4.0/}{Creative Commons Attribution-ShareAlike 4.0}
licence unless stated otherwise. See \href{https://wosm.academy/licence/}{wosm.academy/licence} for details.

\section{Policies}
\begin{itemize}
  \item \textbf{Attendance:} Attendance is not mandatory. The module is self-paced, so you can study the materials when it suits you.
  \item \textbf{Late submissions:} Work submitted after the deadline is penalised.
  \item \textbf{Academic integrity:} All submitted work must be your own.
        See the \href{https://wosm.academy/code-of-conduct/}{Code of Conduct}.
\end{itemize}

\end{document}
